\chapter{Ковариантное отделение переносов вихревой нити}
\label{ch:v6-bosonic-defect-translation-covariant-discretization-gate}

\section{Проблема}

После центрированной релаксации отрицательная мода исчезла, но закрепление
нуля поля в одном узле явно нарушает непрерывную симметрию переносов. Поэтому
малые положительные уровни дискретного гессиана нельзя автоматически считать
физическими массами внутренних колебаний. Требуется отделить положение нити от
её формы.

\section{Поиск решения}

Для стационарного непрерывного профиля введём семейство
\[
  \Phi_{X,Y}(x,y)=\Phi_0(x-X,y-Y).
\]
На однородной плоскости энергия этого семейства не зависит от $X,Y$. Значит,
в координатах $(X,Y,\xi)$, где $\xi$ ортогональны касательным к переносной
орбите, гессиан имеет вид
\[
  \mathcal H=
  \begin{pmatrix}
    0_{2\times2}&0\\
    0&\mathcal H_{\rm int}
  \end{pmatrix}.
\]
Это не новая динамическая гипотеза, а восстановление симметрии, потерянной при
закреплении центра на квадратной сетке. Калибровочно скорректированные
переносные нулевые моды вихря являются стандартным следствием дифференцирования
семейства решений по координатам центра \cite{gustafson-sigal-vortex-stability,
ting-chen-translation-zero-modes}.

Аудит использует три уже стационаризованных центрированных решения. На сетках
$25,33,41$ отрицательных мод нет, а нижние закреплённые уровни равны
соответственно
\[
  0.0120226258,\qquad 0.0008661276,\qquad 0.0004726981.
\]
После добавления $(X,Y)$ два переносных уровня становятся точно нулевыми.

\section{Результат}

\begin{proposition}
Положительная кривизна решёточной энергии по смещению центра не является
массой физической моды. В ковариантном описании она заменяется двумя точными
нулевыми модами переноса. На проверенных центрированных стационарных сетках
оставшийся оператор не имеет отрицательных собственных значений.
\end{proposition}

Однако последний нижний положительный уровень уменьшается при переходе от
сетки $33$ к сетке $41$. Поэтому из этих данных ещё не следует строго
положительный внутренний зазор непрерывного оператора.

\section{Нулевая гипотеза и критерий остановки}

Нулевая гипотеза следующего шага: оставшееся смягчение состоит из угловых
каналов, которые квадратная сетка смешивает и ошибочно расщепляет. Она будет
отвергнута, если после разложения по угловым представлениям обнаружится
сходящаяся отрицательная внутренняя мода. Если разложение не замыкается при
сгущении, декартова схема должна быть заменена полярным оператором, а не
интерпретирована физически.

\section{Статус}

Переносный решёточный артефакт устранён концептуально и вычислительно.
Неотрицательность эффективной подсистемы подтверждена только на проверенных
сетках; непрерывный внутренний зазор, полный гессиан $Q+T+B$, замыкание нити в
хопфову конфигурацию и рождение наблюдаемой материи остаются открытыми.

\begin{thebibliography}{9}
\bibitem{gustafson-sigal-vortex-stability}
S.~Gustafson, I.~M.~Sigal,
\emph{The stability of magnetic vortices},
Commun. Math. Phys. 212 (2000) 257--275, arXiv:math/9904158.

\bibitem{ting-chen-translation-zero-modes}
I.~M.~Sigal, T.~Tzaneteas,
\emph{On stability of Abrikosov vortex lattices},
arXiv:1308.5446.
\end{thebibliography}